\chapter{Literature Review}

\section{Introduction}

Metabolic syndrome (MetS) is a major public health concern characterized by a combination of risk factors such as abdominal obesity, hypertension, insulin resistance, and dyslipidemia, which increases the risk of cardiovascular diseases and type 2 diabetes \citep{peterseim2024metabolic}. Several studies have shown that lifestyle factors-including workplace stress, physical activity, and dietary choices-significantly influence the development and progression of metabolic syndrome \citep{chukwuemeka2025metabolic, gallardo2023diet, chomiuk2024physical}. This paper reviews previous research on these key variables and their interrelationships with metabolic health.

\section{Understanding Metabolic Syndrome}

Metabolic syndrome is not a single disease but a group of metabolic abnormalities that occur together, increasing the risk of chronic illnesses. According to \cite{idf2006_metabolic}, as cited in \cite{alberti2009harmonizing}, an individual is diagnosed with Metabolic Syndrome if they exhibit at least three of the following: abdominal obesity, high blood pressure, high fasting blood sugar, high triglycerides, and low HDL cholesterol levels. \par
The global prevalence of metabolic syndrome has risen significantly in the last two decades, especially in developing countries due to rapid urbanization and lifestyle changes \citep{adeboye2012epidemiology}.

\section{Workplace Stress and Metabolic Syndrome}

Workplace stress is a psychological and physiological response to work-related pressures that exceed an individual's coping ability. Chronic exposure to stress activates the hypothalamic-pituitary-adrenal (HPA) axis, increasing cortisol levels, which can lead to insulin resistance, fat accumulation, and inflammation - all of which are associated with metabolic syndrome \citep{chrousos2009stress}.


\par Studies have found that high job demands, lack of control, and poor social support at work are associated with an increased risk of metabolic abnormalities \citep{griep2015job}. \cite{chandola2006work} found that workers who experienced long-term stress at work had a significantly higher risk of developing Metabolic Syndrome. In the Ghanaian context, workplace stress is often exacerbated by job insecurity, poor work conditions, and long working hours, which can further worsen health outcomes.

\section{Physical Activity and Metabolic Health}

Physical activity plays a vital role in the regulation of metabolic functions. Regular exercise improves insulin sensitivity, reduces abdominal fat, lowers blood pressure, and improves lipid profiles. In contrast, sedentary behavior is associated with an increased risk of metabolic syndrome \citep{seo2023leisure}.

\par \cite{warburton2006health} concluded that moderate to vigorous physical activity significantly reduces the risk of developing metabolic syndrome. However, many employees, particularly in office settings, engage in prolonged sitting during work hours, which contributes to physical inactivity and poor metabolic health. In today's world, physical inactivity is on the rise due to urban lifestyles and transportation habits, especially among professionals \citep{hettiarachchi2021physical}.

\section{Dietary Choices and Metabolic Syndrome}

Diet is another key factor related to the development of metabolic syndrome. Diets high in refined sugars, trans fats, and processed foods contribute to obesity and insulin resistance. However, diets rich in fruits, vegetables, whole grains, and lean proteins have protective effects \citep{siqueira2024major}.


\par \cite{hu2002diet} showed that poor dietary patterns are directly associated with increased risk of metabolic disorders. Also, workplace environments that promote unhealthy eating such as limited access to nutritious food and time constraints that lead to fast food consumption can increase the risk of Metabolic Syndrome among employees. Cultural dietary habits and economic factors also influence food choices.

\section{Interactions Among Stress, Activity, and Diet}

Although these three factors have been studied individually, research indicates that these factors may interact and intensify one another's effects. For instance, individuals under high levels of stress may exercise less and make poorer food choices, which together increase the risk of Metabolic Syndrome. A recent review by \cite{stavitz2025plantbased} underscored the importance of integrated lifestyle interventions—combining diet and structured exercise—to achieve superior metabolic outcomes, including reductions in visceral fat, improved endothelial function, and enhanced glycemic control.


\section{Summary of Gaps in the Literature}

While several international studies have examined workplace stress, physical activity, and dietary habits in relation to metabolic syndrome, limited research has explored their combined impact. There is a need for localized studies that consider cultural, occupational, and economic factors affecting these variables.











